\chapter{Memory Management}
\label{chp:memory}

The \jrtwo\ and \ktwo\ have 2 MB of static RAM, 512 KB of flash memory, and a connector for 256 KB of expansion memory. The 6809 can address only 64 KB. To reach the rest, you use the memory management unit (MMU): it divides the CPU's 64 KB into eight slots of 8 KB and lets you choose which 8 KB block of physical memory appears in each slot.

If you write programs for NitrOS-9, the operating system manages the MMU for you; read section~\ref{sec:mem_nos9} for what that means for your program. If you write raw 6809 programs, read the whole chapter: your program owns the MMU.

\section{Blocks}

The MMU sees physical memory as 256 blocks of 8 KB, numbered \verb+$00+ to \verb+$FF+. A block number is eight bits, and an address within a block is thirteen, so a physical address is 21 bits: the block number times \verb+$2000+, plus the offset.

Most blocks are RAM. Blocks \verb+$40+--\verb+$9F+ depend on the FLASHDIS bit in \verb+MMU_IO_CTRL+ (\verb+$FFA1+): while it is clear, as it is after reset, those blocks are the flash memory and the expansion connector. Set it, and they become RAM too. Blocks \verb+$FC+--\verb+$FF+ are always device memory: VICKY's registers, fonts and color tables, and the text screen. Table~\ref{tab:phys_blocks} shows the layout.

\begin{table}[ht]
    \begin{center}
        \begin{tabular}{|l|l|l|l|} \hline
            Blocks & Physical addresses & FLASHDIS = 0 (reset) & FLASHDIS = 1 \\ \hline\hline
            \verb+$00+--\verb+$3F+ & \verb+$000000+--\verb+$07FFFF+ & RAM (512 KB) & RAM \\ \hline
            \verb+$40+--\verb+$7F+ & \verb+$080000+--\verb+$0FFFFF+ & Flash (512 KB) & RAM \\ \hline
            \verb+$80+--\verb+$9F+ & \verb+$100000+--\verb+$13FFFF+ & Expansion (256 KB) & RAM \\ \hline
            \verb+$A0+--\verb+$FB+ & \verb+$140000+--\verb+$1F7FFF+ & RAM (736 KB) & RAM \\ \hline
            \verb+$FC+--\verb+$FF+ & \verb+$1F8000+--\verb+$1FFFFF+ & Devices & Devices \\ \hline
        \end{tabular}
    \end{center}
    \caption{Physical memory blocks}
    \label{tab:phys_blocks}
\end{table}

With FLASHDIS set you have 2,016 KB of RAM in blocks \verb+$00+--\verb+$FB+; with it clear you have 1,248 KB, plus the flash and the expansion memory. FLASHDIS changes only what the CPU sees. VICKY and the DMA engine use physical addresses directly, without the MMU, and always reach RAM at every block from \verb+$00+ to \verb+$FB+.

\begin{note}
    Older cores have no FLASHDIS bit; their device memory is at blocks \verb+$C0+--\verb+$C3+, and blocks \verb+$A0+--\verb+$FB+ are not all RAM. Bit 7 of \verb+MMU_IO_CTRL+ always reads 1 on a core with FLASHDIS, so your program can test it before setting the bit.
\end{note}

\section{Slots and LUTs}

The CPU's 64 KB is divided into eight slots (table~\ref{tab:mem_slots}). The top three bits of a CPU address select the slot. The MMU looks up the block number for that slot and puts it in front of the remaining 13 bits to make the physical address (figure~\ref{fig:mmu_address_xlate}, where ``bank'' means slot).

\begin{table}[ht]
    \begin{center}
        \begin{tabular}{|c|c|l|} \hline
            Slot & A15--A13 & CPU addresses \\ \hline\hline
            0 & 000 & \verb+$0000+--\verb+$1FFF+ \\ \hline
            1 & 001 & \verb+$2000+--\verb+$3FFF+ \\ \hline
            2 & 010 & \verb+$4000+--\verb+$5FFF+ \\ \hline
            3 & 011 & \verb+$6000+--\verb+$7FFF+ \\ \hline
            4 & 100 & \verb+$8000+--\verb+$9FFF+ \\ \hline
            5 & 101 & \verb+$A000+--\verb+$BFFF+ \\ \hline
            6 & 110 & \verb+$C000+--\verb+$DFFF+ \\ \hline
            7 & 111 & \verb+$E000+--\verb+$FFFF+ \\ \hline
        \end{tabular}
    \end{center}
    \caption{CPU memory slots}
    \label{tab:mem_slots}
\end{table}

\begin{figure}[ht]
    \begin{center}
        \includegraphics[scale=0.65]{images/mmu_address_xlate.pdf}
    \end{center}
    \caption{MMU address translation}
    \label{fig:mmu_address_xlate}
\end{figure}

The eight block numbers for the eight slots form a look-up table (LUT). The MMU holds four LUTs, so you can keep four complete memory maps and switch between them by writing one register. The ACT\_LUT field of \verb+MMU_MEM_CTRL+ (\verb+$FFA0+) selects the LUT that translates addresses, and the EDIT\_LUT field selects the LUT that the slot registers at \verb+$FFA8+--\verb+$FFAF+ read and write. The two fields are independent: you can build a new map in one LUT while another one runs your program. When EDIT\_LUT and ACT\_LUT are the same, a write to a slot register changes the memory map at once.

Every slot can hold any block, including a device block. Slot 7 is different only because part of it is fixed, as the next section describes.

\section{Fixed addresses}

The top three pages of the CPU's address space behave the same in every LUT:

\begin{description}
    \item[\texttt{\$FE00}--\texttt{\$FFFF}] are I/O registers: the interrupt controller, timers, DMA, serial ports, VICKY's control registers and the rest (see the memory map appendix). The MMU registers are at \verb+$FFA0+--\verb+$FFAF+. Reads and writes in this range never reach the RAM block in slot 7, with one exception: \verb+$FFF0+--\verb+$FFFF+.

    \item[\texttt{\$FFF0}--\texttt{\$FFFF}] hold the 6809's interrupt and reset vectors. With VEC\_RAM (bit 1 of \verb+MMU_IO_CTRL+) clear, they come from the block in slot 7, like ordinary memory. With VEC\_RAM set, 16 bytes of RAM inside the FPGA replace them, so the vectors stay put whatever block you map into slot 7.

    \item[\texttt{\$FD00}--\texttt{\$FDFF}] are ordinary memory in slot 7 while CONST\_RAM (bit 0 of \verb+MMU_IO_CTRL+) is clear. Set it, and 256 bytes of RAM inside the FPGA replace them in every LUT. Use this page for code and data that must stay visible whatever the MMU maps: interrupt entry code, or a routine that switches LUTs.
\end{description}

The constant RAM and the vector RAM start with unknown contents. To fill one, set its bit and write to it; with the bit set, your reads and writes go to the internal RAM and the block in slot 7 is untouched.

\section{Device blocks}

Blocks \verb+$FC+--\verb+$FF+ hold the larger VICKY memories (table~\ref{tab:device_blocks}). To reach one, map it into a free slot, read or write it, and map the slot's usual block back. The registers inside each block are described in the chapters on text, graphics, sprites and tiles.

\begin{table}[ht]
    \begin{center}
        \begin{tabular}{|c|l|l|} \hline
            Block & Offsets & Contents \\ \hline\hline
            \verb+$FC+ & \verb+$0000+--\verb+$0BFF+ & Gamma tables (blue, green, red) \\
                       & \verb+$0C00+--\verb+$0FFF+ & Mouse pointer \\
                       & \verb+$1000+--\verb+$10FF+ & Bitmap registers \\
                       & \verb+$1100+--\verb+$11FF+ & Tile registers \\
                       & \verb+$1200+--\verb+$12FF+ & \verify{Memory text mode registers} \\
                       & \verb+$1300+--\verb+$16FF+ & Sprite registers \\
                       & \verb+$1700+--\verb+$173F+ & Text foreground color LUT \\
                       & \verb+$1740+--\verb+$177F+ & Text background color LUT \\ \hline
            \verb+$FD+ & \verb+$0000+--\verb+$0FFF+ & Fonts (two sets of 2 KB) \\
                       & \verb+$1000+--\verb+$1FFF+ & Graphics color LUTs 0--3 (1 KB each) \\ \hline
            \verb+$FE+ & \verb+$0000+--\verb+$12BF+ & Text screen characters (4,800 bytes) \\ \hline
            \verb+$FF+ & \verb+$0000+--\verb+$12BF+ & Text screen colors (4,800 bytes) \\ \hline
        \end{tabular}
    \end{center}
    \caption{Device blocks}
    \label{tab:device_blocks}
\end{table}

\section{The MMU registers}

Table~\ref{tab:mmu_registers} shows the MMU registers, and the list after it describes each one.

% Generated by tools/regs.py from data/registers/mmu.toml. Do not edit.
\begin{table}[ht]
    \begin{center}
        \setlength{\tabcolsep}{3pt}
        \begin{tabular}{|c|l|BBBBBBBB} \hline
            Address & Name & 7 & 6 & 5 & 4 & 3 & 2 & 1 & 0 \\ \hline\hline
            \verb+$FFA0+ & MMU\_\allowbreak{}MEM\_\allowbreak{}CTRL & \bitspan{2}{---} & \bitspan{2}{EDIT\_\allowbreak{}LUT} & \bitspan{2}{---} & \bitspan{2}{ACT\_\allowbreak{}LUT} \\ \hline
            \verb+$FFA1+ & MMU\_\allowbreak{}IO\_\allowbreak{}CTRL & \bitspan{1}{FLASHDIS\_\allowbreak{}OK} & \bitspan{4}{---} & \bitspan{1}{FLASHDIS} & \bitspan{1}{VEC\_\allowbreak{}RAM} & \bitspan{1}{CONST\_\allowbreak{}RAM} \\ \hline
            \verb+$FFA8+--\verb+$FFAF+ & MMU\_\allowbreak{}SLOT\_\allowbreak{}0–7 & \bitspan{8}{Block number} \\ \hline
        \end{tabular}
    \end{center}
    \caption{MMU registers}
    \label{tab:mmu_registers}
\end{table}

% Generated by tools/regs.py from data/registers/mmu.toml. Do not edit.
\begin{description}
    \item[MMU\_MEM\_CTRL] (\verb+$FFA0+, read/write, reset \verb+$00+) Selects the LUT that translates CPU addresses and the LUT that the slot registers edit.
    \begin{description}
        \item[EDIT\_LUT] (bits 5--4) LUT (0–3) that the slot registers at \$FFA8–\$FFAF read and write.
        \item[ACT\_LUT] (bits 1--0) LUT (0–3) that translates CPU addresses.
    \end{description}
    \item[MMU\_IO\_CTRL] (\verb+$FFA1+, read/write, reset \verb+$00+) Controls the constant RAM pages and the flash and expansion windows. Change it with a read-modify-write, so that you keep the bits other code has set.
    \begin{description}
        \item[FLASHDIS\_OK] (bit 7, read-only) Always reads 1 on a core that implements FLASHDIS. Older cores read back the stored 0.
        \item[FLASHDIS] (bit 2) 1 = blocks \$40–\$9F are RAM. 0 = blocks \$40–\$7F are flash and \$80–\$9F are the expansion connector.
        \item[VEC\_RAM] (bit 1) 1 = 16 bytes of internal RAM replace \$FFF0–\$FFFF (the 6809 vectors), in every LUT. 0 = the vectors come from the block mapped in slot 7.
        \item[CONST\_RAM] (bit 0) 1 = 256 bytes of internal RAM replace \$FD00–\$FDFF, in every LUT. 0 = \$FD00–\$FDFF is ordinary mapped memory.
    \end{description}
    \item[MMU\_SLOT\_0] (\verb+$FFA8+--\verb+$FFAF+, read/write) Slot registers 0–7: the block each 8 KB slot maps, in the LUT selected by EDIT\_LUT. At reset LUT 0 maps slots 0–7 to blocks 0–7, except slot 7, which maps flash block \$7F (or RAM block 7 after a debug-port RAM boot). Reset doesn't set LUTs 1–3.
\end{description}


\section{Reset and boot}
\label{pg:mmu_boot_config}

Power-up and the reset button both reset the MMU:

\begin{itemize}
    \item \verb+MMU_MEM_CTRL+ and \verb+MMU_IO_CTRL+ are \verb+$00+: LUT 0 is active and being edited, FLASHDIS is clear, and the constant RAM and vector RAM are off.
    \item LUT 0 maps slots 0--6 to RAM blocks 0--6, and slot 7 to flash block \verb+$7F+, the last 8 KB of flash. The CPU fetches its reset vector from \verb+$FFFE+ in that block, so the machine boots from flash.
    \item LUTs 1--3 keep whatever they held. Set every slot of a LUT before you make it active.
\end{itemize}

For development you can boot from RAM instead. The FoenixMgr command \verb+--boot ram+ sends a command over the USB debug port that makes every later reset of the CPU map RAM block 7 into slot 7 instead of flash block \verb+$7F+. Each FoenixMgr upload resets the CPU this way, so it also starts your program. The setting lasts until you send \verb+--boot flash+, power the machine off, or \verify{press the reset button}. Your program then supplies its own vectors at \verb+$FFF0+--\verb+$FFFF+ of block 7. The raw skeleton in \verb+examples/raw+ works this way; its comments describe the steps.

\section{Raw 6809 examples}

These examples use the names from \verb+examples/include/wildbits_regs.inc+, which the manual's register tables are generated from.

\example{Reaching a device block}

This fragment maps the text color block (\verb+$FF+) into slot 2, fills the screen's colors with white on color 0, and puts RAM block 2 back. It assumes that LUT 0 is both active and being edited, as after reset, and that your program and stack aren't in slot 2.

\begin{verbatim}
            lda   #$FF              text color block
            sta   MMU_SLOT_2        slot 2 ($4000-$5FFF) now shows the colors
            ldx   #$4000
            ldd   #$F0F0            foreground 15, background 0
            ldy   #4800/2           4,800 bytes, two at a time
@fill       std   ,x++
            leay  -1,y
            bne   @fill
            lda   #2                RAM block 2
            sta   MMU_SLOT_2        put it back
\end{verbatim}

\example{Building and switching to a second map}

This fragment fills LUT 1 while LUT 0 keeps running the program, then switches to LUT 1. LUT 1 maps the same blocks as LUT 0 after a RAM boot (0--7), except that slot 3 shows RAM block \verb+$A0+. Keep the slot that holds your running code the same in both LUTs, or the CPU fetches its next instruction from somewhere else.

\begin{verbatim}
            lda   #$10              edit LUT 1, LUT 0 stays active
            sta   MMU_MEM_CTRL
            ldx   #MMU_SLOT_0
            clrb
@slots      stb   ,x+               slot n = block n
            incb
            cmpb  #8
            bne   @slots
            lda   #$A0
            sta   MMU_SLOT_3        slot 3 = block $A0
            lda   #$01              make LUT 1 active (and edit LUT 0)
            sta   MMU_MEM_CTRL
\end{verbatim}

\section{Under NitrOS-9}
\label{sec:mem_nos9}

NitrOS-9 Level 2 takes over the MMU when it boots:

\begin{itemize}
    \item It sets FLASHDIS on a core that has it, so all 2,016 KB of RAM is available to the system and to programs.
    \item It sets CONST\_RAM and VEC\_RAM and keeps its own code in the constant page at \verb+$FD00+ and its vectors in the vector RAM.
    \item LUT 0 holds the system's map. LUT 1 holds the map of the running process; NitrOS-9 reloads it on every task switch.
\end{itemize}

Each process gets its own 64 KB address space, built from the blocks NitrOS-9 allocates to it. \verb+$FD00+--\verb+$FFFF+ belong to the system in every process; the I/O registers are still at \verb+$FE00+--\verb+$FFFF+, but your program can't use that range as memory.

Don't write the MMU registers from a NitrOS-9 program: the next task switch undoes your changes, and a wrong slot crashes the system. Ask NitrOS-9 instead. \verb+F$MapBlk+ maps one or more blocks, such as a device block, into a free part of your address space and returns its address; \verb+F$ClrBlk+ unmaps them. \verb+examples/nos9/bgcolor.asm+ uses the two calls to reach the text color LUTs in block \verb+$FC+:

\begin{verbatim}
            pshs  u                 F$MapBlk returns the address in U
            ldx   #TEXT_LUT_BLK     block $FC (from wildbits.d)
            ldb   #1                one block
            os9   F$MapBlk
            bcs   fail
            stu   <mapaddr          the data area is still reachable through DP
            puls  u
            ...                     use the block through mapaddr
            ldu   <mapaddr
            ldb   #1
            os9   F$ClrBlk          give the address range back
\end{verbatim}

Use the block names that NitrOS-9 defines in \verb+wildbits.d+, such as \verb+VICKY_BLK+ and \verb+COLOR_RAM_BLK+, rather than numbers. To get more memory for data, use \verb+F$Mem+ to grow your data area, or \verb+F$AllRAM+ and \verb+F$MapBlk+ to allocate blocks and map them yourself.
